% !TeX root = ../algebralineare.tex

\section{Applicazioni lineari} \label{sec:linapp}
\subsection{Definizione di applicazione lineare} \label{subsec:linappdef}
    Siano $V, W$ due spazi vettoriali sussistenti su un campo $F$, la mappa $f: V \rightarrow W$ è definita applicazione lineare se soddisfa le seguenti condizioni di linearità:
    \begin{enumerate}[label=(\arabic*)]
        \item $\forall \mathbf{a}, \mathbf{b} \in V, \; f(\mathbf{a} + \mathbf{b}) = f(\mathbf{a}) + f(\mathbf{b})$
        \item $\forall \mathbf{a} \in V, \; \forall \lambda \in F, \; f(\lambda \mathbf{a}) = \lambda f(\mathbf{a})$
    \end{enumerate}

    \vspace{10pt}
    \begin{theorem}[Teorema dell'esistenza ed unicità dell'estensione lineare]
       Siano $V, W$ due spazi vettoriali, sia $B = \{\mathbf{b}_{1}, ..., \mathbf{b}_{n}\}$ una base di $V$ e sia $G$ un insieme di n vettori di $W$, allora esiste una e una sola
       applicazione lineare $T : V \rightarrow W$ tale per cui $T(\mathbf{b}_{i}) = \mathbf{g}_{i}, \; i \in \{1, ..., n\}$.
    \end{theorem}

    \begin{proof}[Dimostrazione]
        Per un generico $\mathbf{v} \in V$, lo si scomponga unicamente secondo $B$ e si definisca un'applicazione $T: V \rightarrow W$ per cui:
        \[
            T(\mathbf{v}) = T(\sum_{i} a_i \mathbf{b}_{i}) = \sum_{i} a_i \mathbf{g}_{i}
        \]
        Dimostrare che $T$ è un'applicazione lineare è triviale.
        \[
            T(\mathbf{v} + \mathbf{u}) = T(\sum_{i} a_i \mathbf{b}_{i} + \sum_{i} c_i \mathbf{b}_{i}) = T(\sum_{i} (a_i + c_i) \mathbf{b}_{i}) = \sum_{i} (a_i + c_i) \mathbf{g}_{i} = \sum_{i} a_i \mathbf{g}_{i} + \sum_{i} c_i \mathbf{g}_{i} = T(\mathbf{v}) + T(\mathbf{u})
        \]
        \[
            T(\lambda \mathbf{v}) = T(\sum_{i} \lambda a_i \mathbf{b}_{i}) = \sum_{i} \lambda a_i \mathbf{g}_{i} = \lambda T(\mathbf{v})
        \]
        Per costruzione, si ha che:
        \[
            T(\mathbf{b}_{i}) = T(1 \cdot \mathbf{b}_{i} + \sum_{j \neq i} 0 \cdot \mathbf{b}_{j}) = \mathbf{g}_{i}
        \]
        Si supponga che vi sia un'altra mappa lineare $L : V \rightarrow W$ tale per cui $L(\mathbf{b}_{i}) = \mathbf{g}_{i}$, allora:
        \[
            L(\mathbf{v}) = L(\sum a_i \mathbf{b}_{i}) = \sum a_i L(\mathbf{b}_{i}) = \sum a_i \mathbf{g}_{i} = T(\mathbf{v})
        \]
    \end{proof}

    \vspace{10pt}
    \begin{theorem}[Teorema dell'immagine di base] \label{thm:image_of_basis}
        Sia $f : V \rightarrow W$ un'applicazione lineare e sia $B = \{\mathbf{b}_{1}, ..., \mathbf{b}_{n}\}$ una base di $V$, allora $f(B) = \{f(\mathbf{b}_{1}), ..., f(\mathbf{b}_{n})\}$ è un generatore di $\textnormal{Im}(f)$.
    \end{theorem}

    \begin{proof}[Dimostrazione]
        Si prenda un $\mathbf{w} \in \textnormal{Im}(f)$, allora $\exists \mathbf{v} \in V$ tale cui $f(\mathbf{v}) = \mathbf{w}$. Dato che $B$ è una base di $V$, si ha:
        \[
            \mathbf{w} = f(\mathbf{v}) = f(\sum_{i} a_i \mathbf{b}_{i}) = \sum_{i} f(a_i \mathbf{b}_{i}) = \sum_{i} a_i f(\mathbf{b}_{i}) \implies \textnormal{Im}(f) \subset \textnormal{span}(f(B))
        \]
        Sia $\mathbf{k} \in \textnormal{span}(\{f(\mathbf{b}_{1}), ..., f(\mathbf{b}_{n})\})$, allora:
        \[
            \mathbf{k} = \sum_{i} c_i f(\mathbf{b}_{i}) = f(\sum_{i} c_i \mathbf{b}_{i}) \implies \exists \mathbf{u} = \sum_{i} c_i \mathbf{b}_{i} \in V \; | \; f(\mathbf{u}) = \mathbf{k} \implies \textnormal{span}(f(B)) \subset \textnormal{Im}(f)
        \]
        Quindi $\textnormal{Im}(f) = \textnormal{span}(f(B))$.
    \end{proof}



\vspace{10pt}
\subsection{Matrici associate e cambi di base} \label{subsec:linappmat}
    Ogni applicazione lineare $f: V \rightarrow W$ può essere definita come:
    \[
        f(\mathbf{v}) = M_{f}^{E, F}\mathbf{v}\rvert_{E}
    \]
    dove $M_{f}^{E, F}$ è la matrice associata ad $f$; $E$ ed $F$ sono, rispettivamente, delle basi di $V$ e $W$ e $\mathbf{v}\rvert_{E}$ è il vettore $\mathbf{v}$ scomposto secondo la base $E$.
    Sia $n$ la dimensione di $V$ e $m$ la dimensione di $W$, allora è possibile definire $M_{f}^{E, F}$ come
    \[
        M_{f}^{E, F} = \begin{bmatrix} f(\mathbf{e}_{1})\rvert_{F} & f(\mathbf{e}_{2})\rvert_{F} & ... & f(\mathbf{e}_{n})\rvert_{F} \end{bmatrix}
    \]

    \vspace{10pt}
    Per operare un cambio di base $\mathbf{v}\rvert_{E} \rightarrow \mathbf{v}\rvert_{E'}$ è necessario determinare la matrice $M_{\textnormal{id}_V}^{E, E'}$ con $E = \{\mathbf{e}_{1}, ..., \mathbf{e}_{n}\}, \; E' = \{\mathbf{e'}_{1}, ..., \mathbf{e'}_{n}\}$ basi di $V$:
    \[
        \mathbf{v}\rvert_{E'} = M_{\textnormal{id}_V}^{E, E'} \mathbf{v}\rvert_{E} = \begin{bmatrix} \mathbf{e}_{1}\rvert_{E'} & \mathbf{e}_{2}\rvert_{E'} & ... & \mathbf{e}_{n}\rvert_{E'} \end{bmatrix} \mathbf{v}\rvert_{E}
    \]
    in altre parole, la matrice associata al cambio di base da $E$ a $E'$ ha per colonne i vettori della base $E$ scritti secondo la base $E'$.
    Inoltre, in virtù del fatto che è possibile passare da una base all'altra in maniera biunivoca, si ha:
    \[
        \mathbf{v}\rvert_{E} =
        M_{\textnormal{id}_V}^{E', E} \mathbf{v}\rvert_{E'} =
        M_{\textnormal{id}_V}^{E', E} (M_{\textnormal{id}_V}^{E, E'} \mathbf{v}\rvert_{E}) \implies
        M_{\textnormal{id}_V}^{E', E} M_{\textnormal{id}_V}^{E, E'} = I \implies
        M_{\textnormal{id}_V}^{E', E} = (M_{\textnormal{id}_V}^{E, E'})^{-1} \; \textnormal{e viceversa}
    \]
    Si indichi d'ora in poi la matrice di cambio base $M_{\textnormal{id}_V}^{E, E'}$ come $P^{E, E'}$

    \vspace{10pt}
    \begin{theorem}[Teorema del cambio di base]
        Siano $V, W$ due spazi vettoriali e siano $E, E'$ basi di $V$, $F, F'$ basi di $W$. Sia $f : V \rightarrow W$ un'applicazione lineare e sia
        $P = P^{E, E'}$, $Q = P^{F, F'}$, $A = M_{f}^{E, F}$, $A' = M_{f}^{E', F'}$ allora $A = Q^{-1} A' P$
    \end{theorem}

    \begin{proof}[Dimostrazione]
        \[
            \begin{tikzcd}[scale cd=1.3, row sep=huge, column sep=huge]
                V_{E'}  \arrow[r, "M_{f}^{E', F'}"]               & W_{F'} \arrow[d, "Q^{-1}"] \\
                V_{E}   \arrow[u, "P"] \arrow[r, "M_{f}^{E, F}"]  & W_{F}
            \end{tikzcd}
        \]
        Come è possibile osservare dal diagramma commutativo soprastante, l'effetto di una mappa da $V_E$ a $W_F$ è identico all'effetto di una composizione delle mappe, rispettivamente,
        da $V_E$ a $V_{E'}$, da $V_{E'}$ a $W_{F'}$, ed infine da $W_{F'}$ a $W_F$. In simboli:
        \[
            \mathbf{u}\rvert_{F} = Q^{-1}\mathbf{u}\rvert_{F'} = Q^{-1} (A' \mathbf{v}\rvert_{E'}) = Q^{-1} A' (P \: \mathbf{v}\rvert_{E}) = (Q^{-1} A' P) \mathbf{v}\rvert_{E} = A \mathbf{v}\rvert_{E}
        \]
    \end{proof}

    \begin{corollary}
        Siano $P, Q \in \mathbb{R}^{n, n}$ e siano $A, A' \in \mathbb{R}^{n, n}$ due matrici tali per cui $A' = Q^{-1}AP$, allora rank($A$) = rank($A'$)
    \end{corollary}



\vspace{10pt}
\subsection{Nucleo di un'applicazione lineare} \label{subsec:linappkernel}
    Sia $f: V \rightarrow W$ un'applicazione lineare, il nucleo (anche chiamato kernel) $\ker(f) := K \subset V$ di $f$ è definito come l'insieme di quei vettori tali per cui $\mathbf{v} \in K \iff f(\mathbf{v}) = \mathbf{0}$.
    E' possibile dimostrare che il kernel di un'applicazione lineare è un sottospazio vettoriale, in quanto:

    \begin{enumerate}[label=(\arabic*)]
        \item $\forall \mathbf{v}, \mathbf{w} \in K, \; f(\mathbf{v} + \mathbf{w}) = f(\mathbf{v}) + f(\mathbf{w}) = \mathbf{0} \implies \mathbf{v} + \mathbf{w} \in K$
        \item $\forall \mathbf{v}, \; \forall \lambda \in \mathbb{R}, \; f(\lambda \mathbf{v}) = \lambda f(\mathbf{v}) = \mathbf{0} \implies \lambda \mathbf{v} \in K$
    \end{enumerate}



\vspace{10pt}
\subsection{Teorema della dimensione} \label{subsec:linappdimtheorem}
    \begin{theorem}[Teorema della dimensione] \label{thm:dimension}
        Sia $f : V \rightarrow W$ un'applicazione lineare fra due spazi vettoriali sussistenti su un campo $F$, allora $\dim(\textnormal{Im}(f)) + \dim(\ker(f)) = \dim(V)$.
    \end{theorem}

    \begin{proof}[Dimostrazione]
        Sia $B = \{ \mathbf{b}_{1}, ..., \mathbf{b}_{n} \}$ una base di $V$. Allora, per il teorema dell'immagine di base [\Cref{thm:image_of_basis}], si ha che $f(B) = \{f(\mathbf{b}_{1}), ..., f(\mathbf{b}_{n})\}$ è un generatore di $\textnormal{Im}(f)$.
        Il teorema di estrazione di base [\Cref{thm:basis_extraction}] garantisce l'esistenza di una base in $f(B)$ e sia $U$ questa base. Sia $B_U = \{\mathbf{b} \in B \; | \; f(\mathbf{b}) \in U\}$ determinato in modo che $|B_U| = |U|$ e sia $B_K = B \setminus B_U$. E' possibile osservare che:
        \[
            V = \mathcal{V}_U \oplus \mathcal{V}_K \; \textnormal{dove} \;\: \mathcal{V}_U = \textnormal{span}(B_U) \; \textnormal{e} \; \mathcal{V}_K = \textnormal{span}(B_K)
        \]
        Dunque è possibile esprimere ogni generico vettore $\mathbf{v} \in V$ come $\mathbf{v} = \mathbf{v}_U + \mathbf{v}_K$ dove $\mathbf{v}_U \in \mathcal{V}_U$ e $\mathbf{v}_K \in \mathcal{V}_K$.
        Non è difficile dimostrare che $\mathcal{V}_U$ è isomorfo a $\textnormal{Im}(f)$ secondo $f$, infatti:
        \[
            \forall \mathbf{v}_U \in \mathcal{V}_U, \; f(\mathbf{v}_U) = f(\sum_{\mathbf{b}_{i} \in B_U} a_i \mathbf{b}_{i}) = \sum_{\mathbf{b}_{i} \in B_U} a_i f(\mathbf{b}_{i})
        \]
        Dato che $U$ è una base di $\textnormal{Im}(f)$, per ogni generico $\mathbf{v}_U \in \mathcal{V}_U$ la scomposizione secondo $U$ di $f(\mathbf{v}_U)$ è unica [\Cref{thm:uniqueness_of_combination}] e coincide con la scomposizione di $\mathbf{v}_U$ secondo $B_U$, innestando una biezione fra $\mathcal{V}_U$ e $\textnormal{Im}(f)$.
        Allora $\dim(\textnormal{Im}(f)) = \dim(\mathcal{V}_U)$.

        \vspace{10pt}
        Si prenda in considerazione il kernel di $f$, ovvero l'insieme di tutti i vettori in $V$ che soddisfano:
        \[
            f(\mathbf{v}) = f(\mathbf{v}_U + \mathbf{v}_K) = \mathbf{0}_W \implies f(\mathbf{v}_U) = -f(\mathbf{v}_K)
        \]

        Anche in questo caso non è difficile trovare una biezione lineare fra $\mathcal{V}_K$ e $\ker(f)$, in quanto:
        \[
            \begin{aligned}
                \forall \mathbf{v}_K \in \mathcal{V}_K \;&\;
                \textnormal{esiste uno e un solo} \;
                \mathbf{v}_U \in \mathcal{V}_U \; | \;
                f(\mathbf{v}_U) = -f(\mathbf{v}_K) \\
                &\implies \textnormal{esiste uno e un solo} \;
                \mathbf{v} = \mathbf{v}_U + \mathbf{v}_K \in \ker(f)
            \end{aligned}
        \]
        Arrivando alla conclusione che $\dim(\ker(f)) = \dim(\mathcal{V}_K)$.

        \vspace{10pt}
        Infine, $\dim(V) = \dim(\mathcal{V}_U) + \dim(\mathcal{V}_K) \implies \dim(V) = \dim(\textnormal{Im}(f)) + \dim(\ker(f))$
    \end{proof}




\vspace{10pt}
\subsection{Endomorfismi lineari semplici e autoaggiunti} \label{subsec:linappend}
    \subsubsection{Endomorfismi lineare semplici} \label{subsubsec:linappsimpleend}
    Un'applicazione lineare $f : V \rightarrow V$ è definita endomorfismo semplice se $V$ è esprimibile come la somma diretta di $k$ sottospazi vettoriali invarianti secondo f e irriducibili ulteriormente. In simboli:
    \[
        V = W_1 \oplus ... \oplus W_k
    \]
    dove ogni $W_i$ è invariante secondo $f$ ($f(W_i) \subset W_i$) e non possiede sottospazi a loro volta invarianti se non $\{\mathbf{0}_V\}$ e $W_i$ stesso.

    \vspace{10pt}
    \subsubsection{Endomorfismi lineari autoaggiunti} \label{subsubsec:linappsadjend}
    Un'applicazione lineare $f: V \rightarrow V$, dove $V$ è uno spazio vettoriale dotato di prodotto interno, è definita endomorfismo autoaggiunto ($self \textnormal{-} adjunct \; endomorphism$) se:
    \[
        \forall \mathbf{v}, \mathbf{w} \in V, \; f(\mathbf{v}) \cdot \mathbf{w} = f(\mathbf{w}) \cdot \mathbf{v}
    \]
    Per determinare se un endomorfismo è autoaggiunto, è necessario dimostrare che per ogni due generici $\mathbf{b}_{1}, \mathbf{b}_{2}$ in una qualsiasi base di $V$ vale
    $f(\mathbf{b}_{1}) \cdot \mathbf{b}_{2} = f(\mathbf{b}_{2}) \cdot \mathbf{b}_{1}$.

    \begin{samepage}
    \begin{theorem}[Teorema degli endomorfismi autoaggiunti] \label{thm:self_adjoint}
        Sia $f: V \rightarrow V$ un endomofismo e sia $E$ una base ortonormale di $V$, allora $f$ è autoaggiunto se e solo se
        $M_{f}^{E, E}$ è simmetrica.
    \end{theorem}

    \begin{proof}[Dimostrazione]
        E' possibile esprimere un generico vettore $\mathbf{v} \in V$ come $\mathbf{v} = (\mathbf{v} \cdot \mathbf{e}_{1})\mathbf{e}_{1} + ... + (\mathbf{v} \cdot \mathbf{e}_{n})\mathbf{e}_{n}$ con $\mathbf{e}_{1}, ..., \mathbf{e}_{n} \in E$. Allora
        \[
            \begin{aligned}
                f(\mathbf{e}_{1}) &= (f(\mathbf{e}_{1}) \cdot \mathbf{e}_{1})\mathbf{e}_{1} + ... + (f(\mathbf{e}_{1}) \cdot \mathbf{e}_{n})\mathbf{e}_{n} \\
                f(\mathbf{e}_{2}) &= (f(\mathbf{e}_{2}) \cdot \mathbf{e}_{1})\mathbf{e}_{1} + ... + (f(\mathbf{e}_{2}) \cdot \mathbf{e}_{n})\mathbf{e}_{n} \\
                \vdots \\
                f(\mathbf{e}_{n}) &= (f(\mathbf{e}_{n}) \cdot \mathbf{e}_{1})\mathbf{e}_{1} + ... + (f(\mathbf{e}_{n}) \cdot \mathbf{e}_{n})\mathbf{e}_{n}
            \end{aligned}
            \; \implies \;
            M_{f}^{E, E} =
            \begin{bmatrix}
                f(\mathbf{e}_{1}) \cdot \mathbf{e}_{1} & f(\mathbf{e}_{2}) \cdot \mathbf{e}_{1} & ... & f(\mathbf{e}_{n}) \cdot \mathbf{e}_{1}   \\
                f(\mathbf{e}_{1}) \cdot \mathbf{e}_{2} & f(\mathbf{e}_{2}) \cdot \mathbf{e}_{2} & ... & f(\mathbf{e}_{n}) \cdot \mathbf{e}_{2}   \\
                \vdots & \ddots & ... & \vdots \\
                f(\mathbf{e}_{1}) \cdot \mathbf{e}_{n} & f(\mathbf{e}_{2}) \cdot \mathbf{e}_{n} & ... & f(\mathbf{e}_{n}) \cdot \mathbf{e}_{n}   \\
            \end{bmatrix}
        \]
        Se $f$ è autoaggiunto, allora ogni coefficiente $f(\mathbf{e}_{i}) \cdot \mathbf{e}_{j} = f(\mathbf{e}_{j}) \cdot \mathbf{e}_{i}$, il che implica che $M_{f}^{E, E}$ è simmetrica.
        Se $M_{f}^{E,E}$ è simmetrica, allora per ogni $\mathbf{e}_{i}, \mathbf{e}_{j} \in E$ si ha che $f(\mathbf{e}_{i}) \cdot \mathbf{e}_{j} = f(\mathbf{e}_{j}) \cdot \mathbf{e}_{i}$, il che implica che $f$ è autoaggiunto.
    \end{proof}
    \end{samepage}




\vspace{10pt}
\subsection{Autovalori, autovettori ed autospazi} \label{subsec:linappeigen}
    Sia $f : V \rightarrow V$ un endomorfismo lineare, si dicono autovettori di $f$ tutti quei vettori $\mathbf{v} \in V$ tali per cui $f(\mathbf{v}) = \lambda \mathbf{v}$, dove $\lambda \in F$, chiamato $autovalore$ di $f$, è uno scalare appartentente al campo su cui $V$ sussiste.

    \vspace{10pt}
    Sia $A = M_{f}^{E, E}$ la matrice associata ad $f$ secondo una base $E$ di $V$, allora gli autovettori sono dati da:
    \[
        f(\mathbf{v}) = A \mathbf{v} = (\lambda I) \mathbf{v} \implies A \mathbf{v} - \lambda I \mathbf{v} = (A - \lambda I) \mathbf{v} = \mathbf{0} \tag{1}
    \]
    Per far sì che vi siano soluzioni non triviali all'equazione (1) il rango di $A - \lambda I$ non deve essere massimo, perchè in quel caso solo il vettore $\mathbf{0}$ sarebbe soluzione.
    In altre parole $f$ possiede autovettori se e solo se:
    \[
        \textnormal{det}(A - \lambda I) = 0 \tag{2}
    \]
    Il termine a sinistra dell'ugualianza in (2) è un polinomio di grado $n$ (dove $n$ è il numero di colonne e righe di $A$) chiamato polinomio caratteristico,
    le cui radici sono i possibili autovalori di $f$. Per il teorema fondamentale dell'algebra il polinomio caratteristico ha $n$ soluzioni, ma spesso (per esempio, quando il campo su cui sussiste $V$ non è algebricamente chiuso) può capitare che
    il numero di autovalori reali sia minore di $n$.

    \vspace{10pt}
    Una volta trovati gli $s$ autovalori distinti $\lambda_1, ..., \lambda_s$ gli autovettori associati sono determinabili come segue:
    \[
        A \mathbf{v} = \lambda \mathbf{v} \implies (A - \lambda I) \mathbf{v} = \mathbf{0} \; \textnormal{con} \; \lambda \in \{\lambda_1, ..., \lambda_s\}
    \]
    L'insieme di tutti gli autovettori associati un autovalore $\lambda$ si dice autospazio ed è definito come $G_{\lambda} = \ker(A - \lambda I)$. $G_{\lambda}$ è sottospazio di $V$ in quanto:
    \begin{enumerate}[label=(\arabic*)]
        \item $\forall \mathbf{v}, \mathbf{w} \in G_{\lambda}, \; f(\mathbf{v} + \mathbf{w}) = \lambda \mathbf{v} + \lambda \mathbf{w} = \lambda (\mathbf{v} + \mathbf{w}) \implies \mathbf{v} + \mathbf{w} \in G_{\lambda}$
        \item $\forall \mathbf{v} \in G_{\lambda}, \forall \mu \in F, \; f(\mu \mathbf{v}) = \mu f(\mathbf{v}) = \mu (\lambda \mathbf{v}) = \lambda (\mu \mathbf{v}) \implies \mu \mathbf{v} \in G_{\lambda}$
    \end{enumerate}

    \vspace{10pt}
    \subsubsection{Molteplicità algebrica e geometrica} \label{subsubsec:linappmultiplicity}
    Per molteplicità algebrica $\mu_f(\lambda)$ di uno specifico autovalore $\lambda$ di $f$ si intende la molteplicità di $\lambda$ come soluzione
    del polinomio caratteristico. In altre parole, se $s$ è il numero di autovalori distinti, allora
    \[
        \textnormal{det}(A - \lambda I) = a \displaystyle \prod_{i = 1}^{s} (\lambda - \lambda_i)^{\mu_f(\lambda_i)}
    \]

    E' opportuno sottolineare che:
    \begin{enumerate}[label=(\arabic*)]
        \item $\forall \lambda \in \{\lambda_1, ..., \lambda_s\}, \; 1 \leq \mu_f(\lambda) \leq n$
        \item $n = \displaystyle \sum_{i}^{s} \mu_f(\lambda_i)$
    \end{enumerate}

    \vspace{10pt}
    Per molteplicità geometrica $\gamma_{f}(\lambda)$ di uno specifico autovalore $\lambda$ di $f$ si intende la dimensione dell'autospazio associato a $\lambda$, ovvero $\gamma_{f}(\lambda) = \dim(G_{\lambda})$.

    \vspace{10pt}
    \begin{osservazione}
        Sia $f: V \rightarrow V$ un endomorfismo lineare, sia $\lambda$ un suo autovalore e sia $\mathbf{v} \in G_{\lambda}$ un autovettore, allora $\lambda \mathbf{v} \in G_{\lambda}$.
    \end{osservazione}

    \vspace{10pt}
    \begin{theorem}[Teorema della scomposizione semplice] \label{thm:simple_decomposition}
        Sia $f : V \rightarrow V$ un endomorfismo lineare e siano $\lambda_1, ..., \lambda_s \in \mathbb{R}$ tutti i possibili autovalori di f. Se $\sum_{i}^{s} \gamma_f(\lambda_i) = n$ allora f è un endomorfismo semplice.
    \end{theorem}

    \begin{proof}[Dimostrazione]
        Dimostrare che $f$ sia un endomorfismo lineare semplice consiste nel dimostrare che è possibile scomporre $V$ in $k$ sottospazi invarianti secondo $f$ e ulteriormente irriducibili. Siano allora gli autospazi $G_{\lambda_i}$ i sopracitati sottospazi.
        Rimane da dimostrare che:
        \begin{enumerate}[label=(\Roman*)]
            \item un generico $G_{\lambda}$ sia invariante secondo $f$ e scomponibile a sua volta in un insieme finito di sottospazi irriducibili.
            \item $\displaystyle \bigoplus_{i}^{s} G_{\lambda_i} = V$
        \end{enumerate}

        \vspace{10pt}
        (I)
        $G_{\lambda}$ è sicuramente invariante secondo $f$, ovvero $f(G_{\lambda}) \subset G_{\lambda}$, in quanto essendo sottospazio di $V$, per chiusura secondo il prodotto scalare-vettore, si ha che
        \[
            \mathbf{v} \in G_{\lambda} \implies \lambda \mathbf{v} = f(\mathbf{v}) \in G_{\lambda}
        \]
        $G_{\lambda}$ è anche pseudo-irriducibile secondo $f$, nel senso che è possibile trovare $p$ sottospazi di $G_{\lambda}$ invarianti e irriducibili. Infatti se $\dim(G_{\lambda}) = 1$ allora gli unici sottospazi ammissibili sono $G_{\lambda}$ stesso e $\{ \mathbf{0}_V \}$ ed è quindi irriducibile.
        Nel caso di $\dim(G_{\lambda}) > 1$ invece si ha:
        \[
            G_{\lambda} = \bigoplus_{j}^{\dim(G_{\lambda})} W_j
        \]
        dove $W_j$ sono i sottospazi monodimensionali di $G_{\lambda}$ tali per cui, data una base $B_{\lambda}$ di $G_{\lambda}$
        \[
            W_j = \textnormal{span}\{ \mathbf{k}_{j} \} \quad \mathbf{k}_{j} \in B_{\lambda}
        \]

        \vspace{10pt}
        (II)
        Per dimostrare che la somma degli autospazi è in primo luogo diretta, è semplicemente necessario dimostrare che, dati due generici autospazi, $G_{\lambda_1} \cap G_{\lambda_2} = \{ \mathbf{0}_V \}$.
        Procedendo per assurdo, sia $\mathbf{v} \in G_{\lambda_1} \cap G_{\lambda_2} \; | \; \mathbf{v} \neq \mathbf{0}_V$. Allora
        \[
           \mathbf{v} \in G_{\lambda_1} \cap G_{\lambda_2} \implies f(\mathbf{v}) = \lambda_1 \mathbf{v} = \lambda_2 \mathbf{v} \implies (\lambda_1 - \lambda_2) \mathbf{v} = \mathbf{0}_V
        \]
        Dato che $\lambda_1 \neq \lambda_2$, in quanto gli autovalori sono esplicitamente distinti, si arriva a $\mathbf{v} = \mathbf{0}_V$, contraddicendo l'ipotesi iniziale.

        \vspace{5pt}
        Infine sia
        \[
            \mathcal{G} = \bigoplus_{i}^{s} G_{\lambda_i}
        \]
        Dato che $\mathcal{G}$ è generato da una somma diretta di sottospazi, è evidente che
        \[
            \mathcal{B} = \bigcup_{i}^{s} B_{\lambda_i}
        \]
        è una base di $\mathcal{G}$, dove $B_{\lambda_i}$ è la base dell'autospazio $G_{\lambda_i}$. La cardinalità di $\mathcal{B}$ è data da
        \[
            |\mathcal{B}| = \sum_{i}^{s} |B_{\lambda_i}| = \sum_{i}^{s} \gamma_{f}(\lambda_i) = n = \dim(V)
        \]
        Quindi $\mathcal{B}$ è un insieme di $n$ vettori linearmente indipendenti (è una base di $\mathcal{G}$) e $\mathcal{B} \subset V$. Allora $\mathcal{B}$ è anche una base di $V$, concludendo che
        \[
            \mathcal{G} = \bigoplus_{i}^{s} G_{\lambda_i} = V
        \]

    \end{proof}




\vspace{10pt}
\subsection{Diagonalizzazione e spettro di una matrice} \label{subsec:linappdiag}
    Sia $A \in \mathbb{R}^{n, n}$ una generica matrice quadrata. Esiste allora un'applicazione lineare $f: V \rightarrow V$ tale per cui $\dim(V) = n$ e $A = M_{f}^{E, E}$ con $E$ base di $V$. $A$ allora si dice diagonalizzabile se:
    \[
        A = PDP^{-1}
    \]
    dove $P$ è una matrice di cambiamento di base dalla ``eigenbase'' $K$ alla base $E$ che ha per colonne gli autovettori che formano una base di $V$ e $D$ è una matrice diagonale sulla quale vi stanno gli autovalori $\lambda_1, ..., \lambda_n$ (non perforza distinti) di $f$
    secondo l'ordine in cui sono messi i vettori in $P$.

    \vspace{10pt}
    Lo spettro di una matrice $A$, oppure, più precisamente, lo spettro di un operatore lineare $f$ su uno spazio vettoriale a dimensione finita $V$ con matrice associata $A = M_{f}^{E, E}$ secondo e verso la base $E$, non è altro che l'insieme di tutti gli autovalori distinti di $f$, spesso
    individuato in simboli come $\Lambda_{A}$ o $\Lambda_{f}$.

    \vspace{10pt}
    E' possibile diagonalizzare una matrice se e solo se $f$ è un endomorfismo lineare semplice su $V$.
    Per computare la diagonalizzazione di una matrice è necessario:
    \begin{enumerate}[label=(\arabic*)]
        \item Individuare lo spettro $\Lambda_{A}$ di $A$ risolvendo $\det(A - \lambda I) = 0$ e facendo in modo che tutti gli autovalori siano reali.
        \item Verificare per ogni autovalore $\lambda \in \Lambda_{A}$ che $\mu_{A}(\lambda) = \gamma_{A}(\lambda)$, potendo così invocare [\Cref{thm:simple_decomposition}].
        \item Trovare una base $K_{\lambda}$ per ogni autospazio $G_{\lambda}$. Dato che $\displaystyle \bigcup_{i}^{s} K_{\lambda_i} = K$ ($K$ è una ``eigenbase'' di $V$), si ha allora $P$.
    \end{enumerate}

    \vspace{10pt}
    \subsubsection{Teorema spettrale} \label{subsubsec:linappspectralthm}
    \begin{theorem}[Teorema spettrale]
        Sia $A = M_{f}^{E, E} \in \mathbb{R}^{n, n}$ una matrice simmetrica associata a $f: V \rightarrow V$ con $E$ base di $V$. Allora $A$ possiede tutti autovalori reali e $V$ è scomponibile secondo $f$ in $s$ autospazi ortogonali.
    \end{theorem}

    \begin{proof}[Dimostrazione]
        (I) Dato che ogni generico autovalore $\lambda$ di $A$ è in $\mathbb{C}$, è necessario dimostrare che $\bar{\lambda} = \lambda$ per assicurare che tutti gli autovalori siano in $\mathbb{R}$.
        Si prenda in considerazione la seguente equazione:
        \[
            (A \mathbf{v})^{T} = (\lambda \mathbf{v})^{T} \implies (A \mathbf{v})^{T} \bar{\mathbf{v}} = (\lambda \mathbf{v})^{T} \bar{\mathbf{v}} \implies \mathbf{v}^{T} A^{T} \bar{\mathbf{v}} = \lambda \mathbf{v}^{T} \bar{\mathbf{v}} \tag{1.1}
        \]
        Posto che $A^{T} = A$ e $\overline{A \mathbf{v}} = \bar{\lambda} \bar{\mathbf{v}} = \bar{A} \bar{\mathbf{v}} = A \bar{\mathbf{v}}$, si ha:
        \[
            \mathbf{v}^{T} (A^{T} \bar{\mathbf{v}}) = \mathbf{v}^{T} (A \bar{\mathbf{v}}) = \mathbf{v}^{T} (\bar{\lambda} \bar{\mathbf{v}}) \implies \bar{\lambda} \mathbf{v}^{T} \bar{\mathbf{v}} = \lambda \mathbf{v}^{T} \bar{\mathbf{v}} \implies \bar{\lambda} = \lambda \in \mathbb{R} \tag{1.2}
        \]
        D'ora in poi denoteremo lo spettro di $f$ come $\Lambda_{f}$.

        \vspace{10pt}
        (II) Per poter affermare che $V$ sia innanzitutto scomponibile secondo i propri autospazi è necessario dimostrare che
        \[
            \sum_{i}^{s} \gamma_f(\lambda_i) = n
        \]
        Inizieremo dimostrando che per ogni $\lambda \in \Lambda_{f}, \; \gamma_{f}(\lambda) = \mu_{f}(\lambda)$. Si assuma per assurdo che vi sia un $\lambda \in \Lambda_{f}$ tale per cui $\gamma_{f}(\lambda) < \mu_{f}(\lambda)$. Per semplicità siano $\gamma = \gamma_{f}(\lambda)$, $\mu = \mu_{f}(\lambda)$. Dato che la molteplicità geometrica di un autovalore corrisponde alla
        dimensione dell'autospazio $G_{\lambda}$ associato a $\lambda$, si consideri una base ortonormale di $G_{\lambda}$ e la si estenda ad una base ortonormale $B$ di $V$. Sia allora $P$ la matrice di cambiamento di base da $B$ a $E$, si ha che
        \[
            A = PDP^{T} \; \textnormal{con} \; D = \begin{bmatrix}
                \lambda I_{\gamma} & K \\
                0 & C
            \end{bmatrix}
        \]
        Per similarità fra $A$ e $D$, $D$ è simmetrica, quindi $K = 0$. Inoltre
        \[
            A = PDP^{T} \implies AP = PD \implies AP - \xi P = PD - \xi P \implies (A - \xi I) P = P (D - \xi I)
        \]
        $A - \xi I$ e $D - \xi I$ sono simili, il che implica che $\det(A - \xi I) = \det(D - \xi I)$. Se i determinanti sono identici, allora anche i polinomi caratteristici lo sono. Sia $\mathcal{C}(\xi)$ il polinomio caratteristico di $A - \xi I$, si ha che
        $\mathcal{C}(\xi)$ ha un fattore $(\xi - \lambda)$ con molteplicità algebrica $\mu$. Si osservi come
        \[
            \mathcal{C}(\xi) = \det(A - \xi I) = \det(D - \xi I) = \det\left(\begin{bmatrix}
                \lambda I_{\gamma} - \xi I_{\gamma} & 0 \\
                0 & C - \xi I_{n - \gamma}
            \end{bmatrix}\right) =
            \det(\lambda I_{\gamma} - \xi I_{\gamma})\det(C - \xi I_{n - \gamma})
        \]
        Dato che $\det(\lambda I_{\gamma} - \xi I_{\gamma})$ è semplicemente $(-1)^{\gamma}(\xi - \lambda)^{\gamma}$, $\det(C - \xi I_{n - \gamma})$ deve contenere un fattore $(\xi - \lambda)^{\mu - \gamma}$ con $\mu - \gamma > 0$.
        Ora, si riprenda la considerazione che $G_{\lambda} = \ker(A - \lambda I)$. Per il teorema della dimensione [\Cref{thm:dimension}] è possibile affermare che
        \[
            \dim(\ker(A - \lambda I)) = n - \rho(A - \lambda I)
        \]
        dove $\rho(A - \lambda I)$ è il rango di $A - \lambda I$. Per similarità fra $A - \xi I$ e $D - \xi I$, per qualsiasi valore di $\xi$, è vero che $\rho(A - \xi I) = \rho(D - \xi I)$. Sia allora $\xi = \lambda$
        \[
            \rho(A - \lambda I) = \rho(D - \lambda I) = \rho\left(\begin{bmatrix}
                0 & 0 \\
                0 & C - \lambda I_{n - \gamma}
            \end{bmatrix}\right) = \rho(C - \lambda I_{n - \gamma})
        \]
        Conseguenza del fatto che $\det(C - \xi I_{n - \gamma})$ abbia un fattore $(\xi - \lambda)$ non nullo, è che $\det(C - \lambda I_{n - \gamma}) = 0$. Il rango di $C - \lambda I_{n-\gamma}$ è minore del rango massimo, che è $n - \gamma$, quindi
        \[
            \rho(C - \lambda I_{n - \gamma}) = \rho(A - \lambda I) < n - \gamma
        \]
        Manipolando un poco l'equazione soprastante, si giunge a
        \[
            n - \rho(A - \lambda I) = \dim(\ker(A - \lambda I)) > \gamma \implies \gamma > \gamma
        \]
        Il che implica che l'ipotesi che $\gamma < \mu$ sia falsa. Sapendo che $\gamma \leq \mu$, si conclude con
        \[
            \forall \lambda \in \Lambda_{f}, \; \gamma = \mu
        \]

        Grazie a [\Cref{thm:simple_decomposition}] si ha la certezza che $f$ sia un endomorfismo semplice secondo i propri autospazi.

        \vspace{10pt}
        (III) Rimane da dimostrare che gli autospazi siano ortogonali fra di loro. Per il teorema degli endomorfismi autoaggiunti [\Cref{thm:self_adjoint}] si ha che $f$ è autoaggiunto. Allora, presi due generici vettori $\mathbf{v} \in G_{\lambda_1}$ e $\mathbf{w} \in G_{\lambda_2}$ appartententi a due autospazi distinti
        \[
            f(\mathbf{v}) \cdot \mathbf{w} = f(\mathbf{w}) \cdot \mathbf{v} \implies \lambda_1 (\mathbf{v} \cdot \mathbf{w}) = \lambda_2 (\mathbf{v} \cdot \mathbf{w}) \implies (\lambda_1 - \lambda_2)(\mathbf{v} \cdot \mathbf{w}) = 0
        \]
        Dato che $\lambda_1$ e $\lambda_2$ sono distinti, $\lambda_1 - \lambda_2 \neq 0$ e quindi $\mathbf{v} \cdot \mathbf{w} = 0$. Tutti gli autospazi sono ortogonali fra di loro.

    \end{proof}

    \vspace{10pt}
    Una conseguenza importante del teorema spettrale è che data una matrice $A$ simmetrica, è sempre possibile trovare una matrice ortogonale che permette di diagonalizzare $A$.
